% 09_conclusion.tex
\section{Conclusion}

We presented Markov Chain Lock (MCL), an active watermark whose
closed-form calibration $S^\star(n_{\min}, \rho, \alpha)$ maps a
regulator-facing specification --- target false-positive rate, text
length, and budget --- directly to the minimum state count, and
whose detector requires only the secret key, no language model
access. The scheme inherits a universal robustness floor
$\delta^\star = 1 - 1/\sqrt{2} \approx 29.3\%$ under the midpoint
threshold convention, and both the calibration and the floor extend
to every $k$-regular transition topology
(\autoref{thm:detection}, \autoref{thm:robustness},
\autoref{thm:k-regular}). Head-to-head against KGW and SWEET on three instruction-tuned LLMs
across four text domains at matched budget, MCL retains substantially
more detection signal under translation and random-substitution
attacks. The empirical post-attack $z$ ratio sits above the
worst-case $(1-\delta_{\mathrm{eff}})^2$ scaling, so we read our
robustness bound as conservative rather than tight
(\autoref{sec:experiments}).
Four directions are immediate: tightening the analytic FPR bound
under non-i.i.d.\ natural-language statistics, where empirical FPR
drifts above the 1\,\% target by 2--7$\times$ at large $S$
(\autoref{tab:cal-sstar}); extending the protocol to longer text
regimes where small $S^\star$ becomes the operative regime;
validating $k$-regular topologies at $k\geq 3$; and characterising
behaviour on base (non-instruction-tuned) LLMs whose entropy
profile differs from the instruct-tuned fleet studied here.
